\documentclass{article}
\usepackage{anyfontsize}
\usepackage[UTF8]{ctex} % 如果需要支持中文
\usepackage{fontspec} 
\setCJKmainfont{SourceHanSansCN-Light}[
    BoldFont=SourceHanSansCN-Medium
]

\usepackage[a4paper, left=0.1in, right=0.1in, top=0.05in, bottom=0.05in]{geometry} % 设置四边页边距为0.1英寸{geometry} % 设置页边距为0.5英寸
\usepackage{enumitem} % 用于自定义列表
\usepackage{amsmath}  % 数学公式支持
\usepackage{hyperref} %不需要目录盒子注释这
\setlength{\parindent}{0pt} % 不缩进
\usepackage{etoolbox} % 用于列表操作
%调整类代码
\usepackage{comment}
\usepackage{tocloft} % 用于定制目录样式
%\usepackage{hyperref} %不需要目录盒子注释这
\usepackage{ifthen}
\usepackage{tikz}
\usepackage{titletoc}
\usepackage{graphicx}
\usepackage{amssymb}  % 提供数学符号，包括\therefore
\usetikzlibrary{positioning}


%目录设定



%快捷键 CMD + / 注释
%  \renewcommand{\cftdot}{} % 隐藏目录中的页码
%  \cftpagenumbersoff{section}
%  \cftpagenumbersoff{subsection}
%  \makeatletter
%  \renewcommand{\@pnumwidth}{0em}  % 设置页码宽度为0
%  \renewcommand{\@tocrmarg}{0em}   % 设置右边距为0
%  \makeatother
 


\usepackage{environ}

\newif\ifshowcontent  % 定义一个条件判断变量
\showcontenttrue    % 设置为 false，隐藏所有 question 环境的内容

\newcounter{question} % 题目计数器
\setcounter{question}{0}
\NewEnviron{question}{%
    \ifshowcontent
        \par\stepcounter{question}\textbf{\thequestion.} \BODY\par % 如果显示内容，则输出
    \fi
}

\NewEnviron{choices}{%
    \ifshowcontent
        \begin{enumerate}[label=\Alph*., leftmargin=*]
            \BODY
        \end{enumerate}\par
    \fi
}

\NewEnviron{correctanswer}{%
    \ifshowcontent
        \textbf{答案:}\BODY\par
    \fi
}



\newenvironment{explanation}
{\textbf{解析:}\par}
{\par}



% 新增考察方面命令
\newcommand{\checklist}[1]{
    \textbf{考察:} #1 \par % 在题目下方显示考察方面
    \addcontentsline{toc}{subsection}{题号\thequestion 考察: #1} % 将考察内容添加到目录
    \vspace{2em} % 添加1em的垂直空间
}

\graphicspath{{images/}} %统一


\begin{document}
 %\fontsize{23}{31}\selectfont
 \tableofcontents % 添加目录
\newpage
%\pagestyle{empty}




\iftrue
\section{模块二MBS考点1:按揭贷款与提前还款}
\setcounter{question}{0}

\begin{question}
    Bennett Bank extends a 5\% APR (annual percentage rate) USD 100,000 30-year mortgage requiring monthly payments. If the mortgage is structured so that it requires interest-only payments for the first 5 years, after which it becomes a self-amortizing mortgage, what would be the portion of the monthly payment applied to the principal in the 61st month?
    \end{question}
    
    \begin{choices}
        \item USD 167.92
        \item USD 174.60
        \item USD 584.59
        \item USD 591.27
    \end{choices}
    
    \begin{correctanswer}
    A
    \end{correctanswer}
    
    \begin{explanation}
    

            \textbf{步骤1：设置贷款参数。}
            \begin{itemize}
                \item \(N = (30 - 5) \times 12 = 300\) （总期数）
                \item \(I/Y = \frac{5}{12} = 0.4167\%\) （月利率）
                \item \(PV = -100,000\) （现值）
                \item \(FV = 0\) （终值）
            \end{itemize}
            
            \textbf{步骤2：计算月供。}
            \begin{itemize}
                \item 使用德州仪器计算器按键：
                \begin{itemize}
                    \item 输入：300 \(\to\) \(N\)
                    \item 输入：0.4167 \(\to\) \(I/Y\)
                    \item 输入：100000 \(\to\) \(PV\)
                    \item 输入：0 \(\to\) \(FV\)
                    \item 按：\(CPT\) \(\to\) \(PMT\)
                \end{itemize}
                \item 计算结果：\(PMT = -584.59\)
            \end{itemize}
            
            \textbf{步骤3：计算第61个月的利息。}
            \begin{itemize}
                \item 利息 = \(100,000 \times \frac{5\%}{12} = 416.67\)
            \end{itemize}
            
            \textbf{步骤4：计算第61个月的本金。}
            \begin{itemize}
                \item 本金 = \(584.59 - 416.67 = 167.92\)
            \end{itemize}
            
            \end{explanation}
            
            \checklist{计算月供和分解利息与本金（使用德州仪器计算器）}


            \begin{question}
                Consider a 30-year mortgage with an initial interest rate of 5\%, compounded monthly, and a principal amount of USD 100,000. After 5 years, the interest rate decreases to 3.5\%. Assuming the refinancing process keeps the principal amount and maturity unchanged, what is the reduction in the monthly repayment amount?
                \end{question}
                
                \begin{choices}
                    \item \$87.78
                    \item \$97.53
                    \item \$112.40
                    \item \$129.67
                \end{choices}
                
                \begin{correctanswer}
               A
                \end{correctanswer}
                
                \begin{explanation}

                    \textbf{步骤1：计算初始月供。}
                    \begin{itemize}
                        \item \(N = 30 \times 12 = 360\)
                        \item \(I/Y = \frac{5\%}{12}\)
                        \item \(PV = 100,000\)
                        \item \(FV = 0\)
                        \item 使用德州仪器计算器按键：
                        \begin{itemize}
                            \item 输入：360 \(\to\) \(N\)
                            \item 输入：0.4167 \(\to\) \(I/Y\)
                            \item 输入：100000 \(\to\) \(PV\)
                            \item 输入：0 \(\to\) \(FV\)
                            \item 按：\(CPT\) \(\to\) \(PMT\)
                        \end{itemize}
                        \item 计算结果：\(PMT = -536.82\)
                    \end{itemize}
                    
                    \textbf{步骤2：计算再融资后的月供。}
                    \begin{itemize}
                        \item \(N = 30 \times 12 = 360\)
                        \item \(I/Y = \frac{3.5\%}{12}\)
                        \item \(PV = 100,000\)
                        \item \(FV = 0\)
                        \item 使用德州仪器计算器按键：
                        \begin{itemize}
                            \item 输入：360 \(\to\) \(N\)
                            \item 输入：0.2917 \(\to\) \(I/Y\)
                            \item 输入：100000 \(\to\) \(PV\)
                            \item 输入：0 \(\to\) \(FV\)
                            \item 按：\(CPT\) \(\to\) \(PMT\)
                        \end{itemize}
                        \item 计算结果：\(PMT = -449.04\)
                    \end{itemize}
                    
                    \textbf{步骤3：计算月供减少额。}
                    \begin{itemize}
                        \item 月供减少额 = \(536.82 - 449.04 = 87.78\)
                    \end{itemize}
                    
                    \end{explanation}
                
                \checklist{计算初始和再融资后月供的步骤。（用德州仪器计算器）}



    \begin{question}
        A fixed-income portfolio manager purchases a seasoned 5.5\% agency mortgage-backed security with a weighted average loan age of 60 months. The current balance on the loans is USD 20 million, and the conditional prepayment rate is assumed to be constant at 0.4\% per year. Which of the following is closest to the expected principal prepayment this month?
        \end{question}
        
        \begin{choices}
            \item USD 1,000
            \item USD 7,000
            \item USD 10,000
            \item USD 70,000
        \end{choices}
        
        \begin{correctanswer}
      B
        \end{correctanswer}

        \begin{explanation}

            \textbf{步骤1：计算月度提前还款率（SMM）。}
            \begin{itemize}
                \item \(\text{SMM} = 1 - (1 - 0.004)^{\frac{1}{12}} = 0.0334\%\)
            \end{itemize}
            
            \textbf{步骤2：计算月供。}
            \begin{itemize}
                \item \(N = 60\)
                \item \(I/Y = \frac{5.5\%}{12}\)
                \item \(PV = -20,000,000\)
                \item \(FV = 0\)
                \item 使用德州仪器计算器按键：
                \begin{itemize}
                    \item 输入：60 \(\to\) \(N\)
                    \item 输入：0.4583 \(\to\) \(I/Y\)
                    \item 输入：20000000 \(\to\) \(PV\)
                    \item 输入：0 \(\to\) \(FV\)
                    \item 按：\(CPT\) \(\to\) \(PMT\)
                \end{itemize}
                \item 计算结果：\(PMT = 382,023.24\)
            \end{itemize}
            
            \textbf{步骤3：计算预期提前还款额。}
            \begin{itemize}
                \item \(\text{PRN} = 290,256.58\)
                \item 预期提前还款额 = \((20,000,000 - 290,356.58) \times 0.0334\% = 6,583.02\)
            \end{itemize}
            
            \end{explanation}

\checklist{计算预期提前还款额的步骤。（用德州仪器计算器）}







\begin{question}
    A portfolio manager at a green energy investment fund is evaluating a pool of sustainable infrastructure bonds. The pool contains three different green bonds financing renewable energy projects. The information on the three bonds is shown as follows:
    
    \begin{center}
    \begin{tabular}{|c|c|c|c|}
    \hline
     & Coupon Rate & Maturity & Principal (USD) \\
    \hline
    Green Bond 1 & 4.5\% & 120 months & 2,500,000 \\
    \hline
    Green Bond 2 & 5.5\% & 180 months & 3,500,000 \\
    \hline
    Green Bond 3 & 6.5\% & 240 months & 4,000,000 \\
    \hline
    \end{tabular}
    \end{center}
    
    What is the correct calculation of the weighted average coupon rate of this green bond pool?
    \end{question}
    
    \begin{choices}
        \item 5.45\%
        \item 5.55\%
        \item 5.65\%
        \item 6.45\%
    \end{choices}
    
    \begin{correctanswer}
    C
    \end{correctanswer}
    
    \begin{explanation}
            \textbf{步骤1：计算投资组合总本金}
            \begin{itemize}
                \item Green Bond 1: \$2,500,000
                \item Green Bond 2: \$3,500,000
                \item Green Bond 3: \$4,000,000
                \item 总本金 = \$2,500,000 + \$3,500,000 + \$4,000,000 = \$10,000,000
            \end{itemize}
            
            \textbf{步骤2：计算每只债券的票息收入}
            \begin{itemize}
                \item Green Bond 1: \$2,500,000 × 4.5\% = \$112,500
                \item Green Bond 2: \$3,500,000 × 5.5\% = \$192,500
                \item Green Bond 3: \$4,000,000 × 6.5\% = \$260,000
                \item 总票息收入 = \$112,500 + \$192,500 + \$260,000 = \$565,000
            \end{itemize}
            
            \textbf{步骤3：计算加权平均票息(WAC)}
            \begin{itemize}
                \item WAC公式：
                \[WAC = \frac{\text{总票息收入}}{\text{总本金}}\]
                \item 代入数值：
                \[WAC = \frac{\$565,000}{\$10,000,000} = 5.65\%\]
            \end{itemize}
            
            \textbf{补充说明：}
            \begin{itemize}
                \item 从这个题可以验证一个特征：由于较大本金的债券票息率较高，因此加权平均值偏向高票息率。
            \end{itemize}
          
    \end{explanation}
    
    \checklist{债券池加权平均票息的计算方法。（记住公式）}





    \begin{question}
        You are analyzing a portfolio of green mortgage-backed securities (Green MBS), where the underlying mortgages are all tied to energy-efficient homes. The pool has a total value of \$500 million with varying prepayment characteristics. Which of the following scenarios would likely result in lower prepayment risk for this Green MBS portfolio?
    \end{question}
        
        \begin{choices}
            \item The government introduces new environmental guarantees that cover both interest and principal payments if borrowers default on their green mortgages.
            
            \item Energy-efficient homes are showing high turnover rates as homeowners frequently upgrade to newer, more efficient properties.
            
            \item Most mortgages in the pool have entered their final years, with relatively small remaining balances and high energy efficiency ratings.
            
            \item Borrowers hold mortgages with 3\% interest rates (including green mortgage discounts), while current market rates have risen to 6\% for similar green mortgages.
        \end{choices}
        
        \begin{correctanswer}
        D
        \end{correctanswer}
        
        \begin{explanation}
        
        \textbf{A选项错误。}
        \begin{itemize}
            \item 政府担保会增加提前还款风险，因为：
            \[\text{违约后回收率} = \text{本金} + \text{未付利息}\]
            这种担保实际上是一种提前还款，增加而非降低了提前还款风险。
        \end{itemize}
        
        \textbf{B选项错误。}
        \begin{itemize}
            \item 高周转率直接影响提前还款概率：
            \[\text{CPR} = 1 - (1 - \text{周转率})^{12}\]
            房产频繁交易会导致提前还款率上升。
        \end{itemize}
        
        \textbf{C选项错误。}
        \begin{itemize}
            \item 贷款后期的提前还款风险反而更高：
            \[\text{月还款额} = \text{本金} \times \frac{r(1+r)^n}{(1+r)^n-1}\]
            其中余额较低时，借款人更容易一次性清偿。
        \end{itemize}
        
        \textbf{D选项正确。}
        \begin{itemize}
            \item 当前市场利率(6\%)远高于持有利率(3\%)时，提前还款动机显著降低：
            \[\text{再融资激励} = \frac{\text{现有利率} - \text{市场利率}}{\text{市场利率}} = \frac{3\% - 6\%}{6\%} = -50\%\]
            负值表明再融资成本过高，借款人更倾向于保持现有低利率贷款。
        \end{itemize}

        \end{explanation}
        
        \checklist{影响MBS提前还款风险的关键因素。（需要知道月还款额公式）}
       




            \begin{question}
                If a pool of mortgage loans begins the month with a balance of \$10,500,000, has a scheduled principal payment of \$54,800, and ends the month with a balance of \$9,800,000, what is the CPR for this month?
            \end{question}
                
                \begin{choices}
                    \item 6.18\%
                    \item 42.24\%
                    \item 53.47\%
                    \item 66.67\%
                \end{choices}
                
                \begin{correctanswer}
                C
                \end{correctanswer}
                
                \begin{explanation}
                    \textbf{步骤1：计算实际提前还款额。}
                    \begin{itemize}
                    \item 实际提前还款 = 实际支付 - 计划支付
                    \item 实际提前还款 = (\$10,500,000 - \$9,800,000) - \$54,800 = \$700,000 - \$54,800 = \$645,200
                \end{itemize}
                
                \textbf{步骤2：计算SMM（月度提前还款率）。}
                \begin{itemize}
                    \item \(\text{SMM} = \frac{\text{实际提前还款}}{\text{期初余额} - \text{计划支付}}\)
                    \item \(\text{SMM} = \frac{\$645,200}{\$10,500,000 - \$54,800} = 0.06177\)
                \end{itemize}
                
                \textbf{步骤3：计算CPR（年度提前还款率）。}
                \begin{itemize}
                    \item \(\text{CPR} = 1 - (1 - \text{SMM})^{12}\)
                    \item \(\text{CPR} = 1 - (1 - 0.06177)^{12} = 0.5347 = 53.47\%\)
                \end{itemize}
                    
                    \end{explanation}
                
                \checklist{如何计算按揭贷款池的CPR（年度提前还款率）}




    \begin{question}
        An investor is evaluating the risks associated with Mortgage-Backed Securities (MBS), which are backed by a pool of insured mortgages. These mortgages carry prepayment risk, allowing borrowers to repay early due to favorable interest rate changes. From an investor's perspective, holding an MBS is similar to having a long position in a non-prepayable mortgage pool combined with which of the following options?
        \end{question}
        
        \begin{choices}
            \item A long American put option on the underlying pool of mortgages.
            \item A short American call option on the underlying pool of mortgages.
            \item A long European call option on the underlying pool of mortgages.
            \item A short European put option on the underlying pool of mortgages.
        \end{choices}
        
        \begin{correctanswer}
        B
        \end{correctanswer}
        
        \begin{explanation}
                \textbf{B选项正确。}
                \begin{itemize}
                    \item MBS的提前还款风险类似于持有短期看涨期权，因为借款人有权在利率下降时提前还款。
                    \item 这种提前还款的选择权对投资者来说是一种风险，因为它限制了MBS的价格升值潜力。
                    \item 当利率下降时，借款人更可能提前还款，导致投资者无法从利率下降中获得全部收益。
                    \item 这种情况类似于卖出看涨期权，投资者面临被行权的风险，限制了潜在收益。
                \end{itemize}
        
        \end{explanation}
        
        \checklist{分析MBS风险特征。（可以类比期权特性）}



        \begin{question}
            A homeowner has a 30-year, 4.5\% fixed-rate mortgage with a current balance of USD 300,000. Mortgage rates have been decreasing. If the existing mortgage was refinanced into a new 30-year, 3.5\% fixed-rate mortgage, which of the following is closest to the amount that the homeowner would save in monthly mortgage payments?
            \end{question}
            
            \begin{choices}
                \item USD 170
                \item USD 180
                \item USD 190
                \item USD 200
            \end{choices}
            
            \begin{correctanswer}
            A
            \end{correctanswer}
            
            \begin{explanation}
            
            \textbf{步骤1：计算当前月供。}
            \begin{itemize}
                \item 使用4.5\%利率计算。
                \item \(N = 360\)（30年）
                \item \(I/Y = \frac{4.5}{12}\)
                \item \(PV = 300,000\)
                \item 使用计算器计算月供，结果为USD 1,520。
            \end{itemize}
            
            \textbf{步骤2：计算再融资后的月供。}
            \begin{itemize}
                \item 使用3.5\%利率计算。
                \item \(N = 360\)（30年）
                \item \(I/Y = \frac{3.5}{12}\)
                \item \(PV = 300,000\)
                \item 使用计算器计算月供，结果为USD 1,347。
            \end{itemize}
            
            \textbf{步骤3：计算月供节省额。}
            \begin{itemize}
                \item 月供节省额 = 当前月供 - 再融资后月供。
                \item 月供节省额 = USD 1,520 - USD 1,347 = USD 173。
            \end{itemize}
            
            \end{explanation}
            
            \checklist{计算月供和节省额的步骤（用德州仪器计算器）}



            \begin{question}
                A financial advisor is reviewing the prepayment penalty policies and types of loans available to clients. Based on the descriptions of prepayment penalties and loan types, which of the following statements is correct?
            \end{question}
            
            \begin{choices}
                \item The interest rate on Adjustable Rate Mortgages (ARMs) is fixed, typically based on LIBOR + 250bps.
                \item Mortgage interest rates are always variable, allowing borrowers to prepay when rates drop to lower their interest rate.
                \item Prepayment penalty policies are the same across all banks, typically not charging penalties within the first two years of the loan.
                \item The prepayment option allows borrowers to refinance at a lower rate when interest rates drop.
            \end{choices}
            
            \begin{correctanswer}
                D
            \end{correctanswer}

            \begin{explanation}
                \textbf{A选项错误。}
                \begin{itemize}
                    \item 可变利率房贷（ARMs）的利率是浮动的，通常基于“基准利率 + 利差”。例如，LIBOR + 250bps。这意味着利率会随着市场利率的变化而变化，而不是固定的。
                \end{itemize}
            
                \textbf{B选项错误。}
                \begin{itemize}
                    \item 固定利率房贷的利率是固定的，在整个贷款期限内不会改变。并非所有房贷利率都是浮动的。
                \end{itemize}
            
                \textbf{C选项错误。}
                \begin{itemize}
                    \item 提前还款罚息政策在不同的银行有所不同，通常与提前还款的时间段相关。例如，借款人在贷款后1年内提前还款，征收2\%的罚息；借款人在贷款后1年至2年内提前还款，征收1\%的罚息；借款人在贷款后2年以上提前还款，不征收罚息。
                \end{itemize}
            
                \textbf{D选项正确。}
                \begin{itemize}
                    \item 提前还款选项允许借款人在利率下降时提前还款，从而以较低的利率重新融资。
                \end{itemize}
            \end{explanation}

\checklist{提前还款罚息政策和贷款类型的区别。(提前还款罚息政策在不同的银行有所不同)}  


\begin{question}
    A junior risk analyst at an asset management firm is tasked with evaluating different types of loans and their associated risks for a client portfolio. Based on the descriptions of loan types and their risks, which of the following descriptions is correct?
\end{question}

\begin{choices}
    \item Standard loans are issued by private institutions and typically do not meet loan limit and minimum down payment requirements.
    \item Mortgage-Backed Securities (MBS) are riskier because the underlying loans have a higher default risk.
    \item Non-conforming loans are recognized by the three major mortgage agencies and are typically provided by government-backed institutions.
    \item Non-conforming loans are usually issued by private institutions outside the three major agencies and may carry higher risks.
\end{choices}

\begin{correctanswer}
    D
\end{correctanswer}

\begin{explanation}

    \textbf{A选项错误。}
    \begin{itemize}
        \item 标准贷款由政府支持的机构提供，满足贷款额度限制和最低首付比例的要求。它们通常由房利美、房地美和吉利美等三大房贷机构认可。
    \end{itemize}
    
    \textbf{B选项错误。}
    \begin{itemize}
        \item 抵押贷款支持证券（MBS）通常由符合标准的贷款构成，其基础资产的贷款违约风险较低，因此MBS的风险较低。
    \end{itemize}
    
    \textbf{C选项错误。}
    \begin{itemize}
        \item 非标准贷款不满足相关标准，通常不被三大房贷机构认可，而是由私人机构发行。
    \end{itemize}
    
    \textbf{D选项正确。}
    \begin{itemize}
        \item 非标准贷款通常由三大机构以外的私人机构发行，可能具有更高的风险。这些贷款可能超出额度限制，或在某些方面未能符合标准。
    \end{itemize}

\end{explanation}

\checklist{标准贷款和非标准贷款的区别。(MBS风险较低)}



\begin{question}
    A loan officer at a bank is reviewing a client's mortgage repayment schedule. The client has an outstanding principal balance of ¥150,000 at the end of the month. The scheduled repayment for this month is ¥10,000, but the client actually repaid ¥15,000. Calculate the Single Monthly Mortality (SMM) rate for this month.
\end{question}

\begin{choices}
    \item 33.33\%
    \item 50\%
    \item 10\%
    \item 20\%
\end{choices}

\begin{correctanswer}
    A
\end{correctanswer}

\begin{explanation}

\textbf{单月提前还款率（SMM）的定义：}
\begin{itemize}
    \item 单月提前还款率（SMM）是指每月提前还款金额占当月未偿还本金余额的比重，其具体公式如下：
    \[
    \text{SMM} = \frac{\text{当月提前还款金额}}{\text{本月底未偿还本金} - \text{本月按期应还本金}}
    \]
\end{itemize}

\textbf{计算过程：}
\begin{itemize}
    \item 本月未偿还本金余额 = ¥150,000
    \item 本月按期应还本金 = ¥10,000
    \item 提前还款金额 = ¥15,000
\end{itemize}

根据公式：
\[
\text{SMM} = \frac{15,000}{150,000 - 10,000} = \frac{15,000}{140,000} = 0.1071
\]

\end{explanation}

\checklist{单月提前还款率（SMM）的计算。（记住公式即可）}



\begin{question}
    The PSA (Public Securities Association) prepayment benchmark assumes that the monthly prepayment rate for a mortgage pool increases with its age or seasonal changes. The 100\% PSA (or simply 100 PSA) for a 30-year mortgage assumes the following increasing CPRs (Conditional Prepayment Rates): CPR starts at 0.2\% in the first month and increases by 0.2\% each month until the 30th month. For example, the CPR for the 14th month is 14(0.2\%) = 2.8\%. From the 30th month to the 360th month, the CPR is 6\%. Based on this information, which of the following statements is correct?
\end{question}

\begin{choices}
    \item 100\% PSA assumes a CPR of 6\% for all months.
    \item 50\% PSA refers to half of the CPR specified by 100\% PSA, while 200\% PSA refers to twice the CPR specified by 100\% PSA.
    \item The PSA benchmark does not account for seasonal variations in the mortgage pool.
    \item The CPR for the 14th month is 2.8\%, which is consistent with the 100\% PSA specification.
\end{choices}

\begin{correctanswer}
    D
\end{correctanswer}

\begin{explanation}

    \textbf{A选项错误。}
    \begin{itemize}
        \item 100\% PSA并不假设所有月份的CPR均为6\%。在前30个月，CPR从0.2\%开始，每月增加0.2\%，直到达到6\%后保持不变。
    \end{itemize}
    
    \textbf{B选项错误。}
    \begin{itemize}
        \item 虽然50\% PSA和200\% PSA确实是100\% PSA的CPR的一半和两倍，但题目要求的是对PSA模型的具体描述，而不是比例关系。
    \end{itemize}
    
    \textbf{C选项错误。}
    \begin{itemize}
        \item PSA预付款基准确实考虑了抵押贷款池的季节性变化和年龄增长，因为它假设月度预付款率会随着时间的推移而增加。
    \end{itemize}
    
    \textbf{D选项正确。}
    \begin{itemize}
        \item 根据100\% PSA的定义，第14个月的CPR计算为\(14 \times 0.2\% = 2.8\%\)。这与PSA模型的规定一致，说明该选项正确。
    \end{itemize}
    
    \end{explanation}

\checklist{PSA预付款基准。（这个概念需要理解）}





\fi


\iftrue
\section{MBS考点:Classifications of MBSs}
\setcounter{question}{0}



\begin{question}
    Tom is an experienced fixed-income trader, and he's considering two different financing strategies by managing his portfolio: a repurchase agreement (repo) and a dollar roll. The portfolio includes an MBS with a par value of USD 100 million with a 30-year FNMA pool with a 4.5\% coupon. He wants to understand the key differences between these strategies. Which of the following statements is true?
\end{question}
    
    \begin{choices}
        \item In a repo, Tom would agree to buy back the same securities he initially sells, while in a dollar roll, he may buy similar but not the same securities.
        \item Interest is added to the price at which the securities are repurchased for both repo and dollar roll transactions.
        \item When Tom evaluates the value of the dollar roll trade, the cost should include the interest earned on the proceeds of the sale for the financing period.
        \item In the dollar roll, Tom should have the coupon and the principal repayment that would have been received on the pool sold during the financing period.
    \end{choices}
    
    \begin{correctanswer}
    A
    \end{correctanswer}
    
    \begin{explanation}

        \textbf{A选项正确。}
        \begin{itemize}
            \item 在回购交易中，Tom会同意在未来以相同价格买回他最初卖出的相同证券。这种交易提供了确定性和安全性。
            \item 在美元滚动交易中，Tom可能会买回类似但不完全相同的证券，这提供了灵活性，但也可能引入不同的提前还款风险。
        \end{itemize}
        
        \textbf{B选项错误。}
        \begin{itemize}
            \item 回购和美元滚动交易的价格中不包含利息。回购价格是事先确定的，而美元滚动交易的价格不包括利息调整。
        \end{itemize}
        
        \textbf{C选项错误。}
        \begin{itemize}
            \item 在美元滚动交易中，Tom出售MBS后，获得的资金不包括利息收入。利息收入由买方在持有期间获得。
        \end{itemize}
        
        \textbf{D选项错误。}
        \begin{itemize}
            \item 在美元滚动交易中，Tom出售MBS后，利息和本金的收入不再属于Tom，而是属于购买这些证券的对手方。
        \end{itemize}
        
        \end{explanation}
    
    \checklist{分析回购和美元滚动交易的关键区别。（美元滚动交易买的是TBA）}



\begin{question}
    A junior trader on a derivatives trading desk is analyzing the differences between dollar roll transactions and traditional repurchase agreements to optimize the firm's trading strategies. What is a key difference between a dollar roll transaction and a traditional repurchase agreement (repo)?
\end{question}
    
    \begin{choices}
        \item In a dollar roll, the initiating party sells securities and agrees to buy them back at a higher price, similar to a repo.
        \item Dollar rolls involve the exchange of securities with the same characteristics between parties.
        \item Dollar rolls may involve receiving different securities in the second month, unlike repos where the same securities are repurchased.
        \item Dollar roll transactions add interest to the repurchase price, unlike repos where interest is not considered.
    \end{choices}
    
    \begin{correctanswer}
    C
    \end{correctanswer}
    
    \begin{explanation}

        \textbf{A选项错误。}
        \begin{itemize}
            \item 在回购协议中，证券以预定价格买回，通常相同或略高于卖出价，但美元回购不一定遵循这一模式。美元回购的价格可能不包括利息调整，重点在于证券的流动性管理。
        \end{itemize}
        
        \textbf{B选项错误。}
        \begin{itemize}
            \item 美元回购允许在第二个月收到不同的证券，这与回购协议中必须买回相同证券的要求不同。这种灵活性使得美元回购在管理不同证券的流动性和风险时更具优势。
        \end{itemize}
        
        \textbf{C选项正确。}
        \begin{itemize}
            \item 美元回购可能涉及不同的证券，这提供了灵活性，但也可能引入不同的提前还款风险和收益特征。这种差异使得美元回购在市场波动时具有独特的风险管理优势。
        \end{itemize}
        
        \textbf{D选项错误。}
        \begin{itemize}
            \item 美元回购不在回购价格中添加利息，发起方在持有期间失去利息收入，而回购协议通常不涉及此类利息损失。这意味着美元回购的收益计算需要考虑利息损失的影响。
        \end{itemize}
        
        \end{explanation}
    
    \checklist{分析美元回购和传统回购协议的关键区别。（美元回购可能涉及不同的证券）}



    \begin{question}
        A financial analyst is evaluating different types of Mortgage-Backed Securities (MBS) for a client. MBS can be classified into Agency MBS, issued by government agencies or sponsored enterprises, and Non-Agency MBS, issued by private entities. Which of the following descriptions is correct?
    \end{question}
    
    \begin{choices}
        \item Agency MBS and Non-Agency MBS are both issued by government agencies.
        \item Agency MBS provides investors with protection against default risk, but not prepayment.
        \item Non-Agency MBS is issued by government-sponsored enterprises, providing investors with protection against default risk.
        \item MBS issued by Ginnie Mae (GNMA) has no implicit guarantee.
    \end{choices}
    
    \begin{correctanswer}
        B
    \end{correctanswer}
    
    \begin{explanation}
    
    \textbf{A选项错误。}
    \begin{itemize}
        \item 机构MBS由政府机构或政府赞助企业发行，而非机构MBS由私人实体发行。两者的发行主体不同。
    \end{itemize}
    
    \textbf{B选项正确。}
    \begin{itemize}
        \item 机构MBS为投资者提供违约保护，但不包括提前还款的保护。这是因为政府机构或赞助企业提供了违约担保。所以该投资者不会面临违约风险，但是会提前还款风险。
    \end{itemize}
    
    \textbf{C选项错误。}
    \begin{itemize}
        \item 非机构MBS由私人实体发行，不提供违约风险的保护，因为没有政府或赞助企业的担保。
    \end{itemize}
    
    \textbf{D选项错误。}
    \begin{itemize}
        \item 吉利美（Ginnie Mae）发行的MBS虽然没有明确的政府担保，但通常被认为有隐性担保，提供一定的安全性。
    \end{itemize}
    
    \end{explanation}
\checklist{MBS的分类。（重点是机构MBS和非机构MBS的区别。）}



\begin{question}
    An investment analyst is assessing the characteristics of Mortgage Pool Securities (MPS) within Agency Mortgage-Backed Securities (MBS) for a client portfolio. Based on the definition and characteristics of MPS, which of the following descriptions is correct?
\end{question}

\begin{choices}
    \item MPS allows investors to receive different returns based on their investment amounts.
    \item MPS involves no fees, and investors receive the entire cash flow from the mortgage pool.
    \item MPS investors must bear prepayment risk, which is the uncertainty of the cash flow amounts they receive.
    \item MPS is issued by private entities and does not provide protection against default risk for investors.
\end{choices}

\begin{correctanswer}
    C
\end{correctanswer}

\begin{explanation}
    \textbf{A选项错误。}
    \begin{itemize}
        \item MPS中所有投资者获得相同的回报，因为现金流是按比例分配的，而不是根据投资份额的不同而变化。
    \end{itemize}

    \textbf{B选项错误。}
    \begin{itemize}
        \item MPS涉及费用，这些费用用于担保和服务抵押贷款，因此投资者不会获得抵押贷款池的全部现金流。
    \end{itemize}

    \textbf{C选项正确。}
    \begin{itemize}
        \item MPS投资者必须接受提前还款风险，这意味着他们收到的现金流金额可能会因为借款人提前还款而减少，导致现金流的不确定性。
    \end{itemize}

    \textbf{D选项错误。}
    \begin{itemize}
        \item MPS是由政府机构或政府赞助企业发行的，而不是私人实体，并且通常提供投资者免受违约风险的保护。
    \end{itemize}
\end{explanation}

\checklist{MPS的特点（MPS的特点包括相同回报、费用扣除和提前还款风险。）}



\begin{question}
    A junior mortgage analyst at a fixed-income investment firm is currently assessing the trading strategies for mortgage-backed securities. The analyst is particularly focused on understanding the TBA (To-Be-Announced) trading mode to optimize client portfolios. Which of the following descriptions is correct regarding the TBA trading mode?
\end{question}

\begin{choices}
    \item In a specified pool trading mode, investors only learn the specific characteristics of the asset pool after the transaction is completed, increasing transaction uncertainty.
    \item In the TBA trading mode, both parties know the specific asset pool at the time of transaction completion, reducing transaction uncertainty.
    \item In the TBA trading mode, the contract requires the buyer to choose the cheapest deliverable asset pool that meets the contract requirements.
    \item The specified pool trading mode provides higher transaction flexibility and liquidity compared to the TBA mode, as all characteristics are known before the transaction.
\end{choices}

\begin{correctanswer}
    C
\end{correctanswer}

\begin{explanation}
    \textbf{A选项错误。}
    \begin{itemize}
        \item 在特定资产池交易模式下，资产池的特征在交易前就已经被清晰定义，这减少了交易的不确定性。
    \end{itemize}

    \textbf{B选项错误。}
    \begin{itemize}
        \item 在TBA模式中，交易的资产池在交易达成时是一类资产类型，而不是某个特定的资产集合，这增加了交易的灵活性。
    \end{itemize}

    \textbf{C选项正确。}
    \begin{itemize}
        \item 在TBA模式下，合约确实要求买方选择一个最便宜可交割的资产池，即在满足合约基本要求的前提下，选择成本最低的资产池进行交割。
    \end{itemize}

    \textbf{D选项错误。}
    \begin{itemize}
        \item 虽然特定资产池模式在交易前就明确了资产池的所有具体特征，但TBA模式由于只确认资产池的类别而不具体限定某一个特定标的，因此具有更高的交易灵活性和流动性。
    \end{itemize}
\end{explanation}
\checklist{TBA的特点。（TBA的特点包括交易灵活性和流动性。）}



\begin{question}
    A financial advisor at an investment bank is analyzing the mechanics of dollar roll transactions in mortgage-backed securities (MBS) trading to provide insights for clients. Based on the detailed description of dollar roll transactions, which of the following is the most accurate description?
\end{question}

\begin{choices}
    \item In a dollar roll transaction, investors must purchase identical securities in two consecutive months to ensure consistency and predictability.
    \item During a dollar roll transaction, the repurchase price in the second month typically includes accrued interest to reflect the time value of money and opportunity cost.
    \item In a dollar roll transaction, the securities purchased next month may differ from the first month, and there is no interest added to the repurchase price, causing one party to lose interest income for a month while the other gains it.
    \item Dollar roll transactions are identical to standard repurchase agreements in all aspects, applicable to all types of MBS trading, and do not affect investors' cash flow.
\end{choices}

\begin{correctanswer}
    C
\end{correctanswer}

\begin{explanation}
    \textbf{A选项错误。}
    \begin{itemize}
        \item 美元滚动交易并不要求投资者在连续两个月内购买完全相同的证券，实际上，下个月购买的证券可能与第一个月不同。
    \end{itemize}

    \textbf{B选项错误。}
    \begin{itemize}
        \item 在美元滚动交易中，第二个月的回购价格上没有增加利息，这与标准的回购协议不同，后者通常会包括应计利息。
    \end{itemize}

    \textbf{C选项正确。}
    \begin{itemize}
        \item 美元滚动交易中，下个月购买的证券可能与第一个月不同，且在回购价格上没有利息的增加，这导致交易的一方在一个月内失去利息收入，而另一方则获得相应的利息收益。
    \end{itemize}

    \textbf{D选项错误。}
    \begin{itemize}
        \item 美元滚动交易与标准的回购协议在关键方面存在区别，特别是在利息处理和证券选择上，因此它们并不完全相同，也不适用于所有类型的MBS交易。
    \end{itemize}
\end{explanation}

\checklist{美元滚动交易的特点。（美元滚动交易的特点包括证券选择、利息处理和交易灵活性。）}





\begin{question}
    Suppose a 4\% pool with a face value of \$2,000,000 was sold in August for 104.00 and repurchased in September for 103.25. The values 104 and 103.25 are percentages of the face value. The monthly payment date is the 15th day of each month(assuming 30 days per month).Funds from the first month's sale will earn one month's interest at 0.15\%. Assume that the forgone coupon and principal payments are 0.6\% of the face value. Calculate the value of this dollar roll transaction.
\end{question}

\begin{choices}
    \item \$5,125
    \item \$6,125
    \item \$7,125
    \item \$8,125
\end{choices}
    \begin{correctanswer}
    B
    \end{correctanswer}
    \begin{explanation}
    \textbf{理论公式}
    \begin{itemize}
    \item value of a dollar roll=A-B+C-D = A - B + C - D
    \item A = 第一个月卖出池子的价格，包括应计利息
    \item B = 第二个月买入池子的价格，包括应计利息
    \item C = 第一个月销售所得资金获得的一个月利息
    \item D = 第一个月卖出池子时放弃的票息和本金支付
    \end{itemize}


    \textbf{步骤1：计算A的值（8月份售出价格加上应计利息）}
    \begin{itemize}
        \item 售出价格 = 200万美元 × 1.04 = 208万美元
        \item 应计利息 = 15/30 × 0.04/12 × 200万美元 = 3333美元
        \item A = 208万美元 + 3333美元 = 2,083,333美元
    \end{itemize}
    
    \textbf{步骤2：计算B的值（9月份回购价格加上应计利息）}
    \begin{itemize}
        \item 回购价格 = 200万美元 × 1.0325 = 206.5万美元
        \item 应计利息 = 15/30 × 0.04/12 × 200万美元 = 3333美元
        \item B = 206.5万美元 + 3333美元 = 2,068,333美元
    \end{itemize}
    
    \textbf{步骤3：计算C的值（第一个月的销售所得资金按0.15\%的利率获得的一个月利息）}
    \begin{itemize}
        \item C = 2,083,333美元 × 0.0015 = 3,125美元
    \end{itemize}
    
    \textbf{步骤4：计算D的值（放弃的票息和本金支付）}
    \begin{itemize}
        \item D = 200万美元 × 0.006 = 12,000美元
    \end{itemize}
    
    \textbf{步骤5：计算美元滚动交易的价值}
    \begin{itemize}
        \item 价值 = A - B + C - D
        \item 价值 = 2,083,333美元 - 2,068,333美元 + 3,125美元 - 12,000美元 = 6,125美元
    \end{itemize}

\end{explanation}

\checklist{美元滚动交易的价值计算。（记住公式即可）}



\begin{question}
    In the modern CMO (Collateralized Mortgage Obligation) structure, how does the relationship between PAC (Planned Amortization Class) bonds and support tranches affect their respective prepayment risks?
\end{question}

\begin{choices}
    \item The PAC tranches do not have a priority claim against the cash flows generated by the underlying mortgage pool.
    \item The PAC tranches have a priority claim against the cash flows generated by the underlying mortgage pool.
    \item The prepayment risk of the PAC tranche is directly proportional to the prepayment risk of the support tranche.
    \item The prepayment risks of the PAC tranche and the support tranche are independent and do not affect each other.
\end{choices}

\begin{correctanswer}
    C
\end{correctanswer}

\begin{explanation}
    \textbf{A选项错误。}
    \begin{itemize}
        \item PAC层的提前还款风险与支持层的提前还款风险之间存在反比关系。也就是说，当PAC层的提前还款风险增加时，支持层的风险会相应减少，反之亦然。
    \end{itemize}

    \textbf{B选项错误。}
    \begin{itemize}
        \item 虽然支持层确实为PAC层提供了一定的提前还款保护，但这并不意味着PAC层在所有情况下都能按时收到预定的本金，特别是在市场出现异常提前还款的情况下，PAC层的收益可能会受到影响。
    \end{itemize}

    \textbf{C选项正确。}
    \begin{itemize}
        \item PAC层的提前还款风险减少确实会导致支持层的提前还款风险增加。这是因为PAC层的稳定性会影响到支持层的现金流，二者之间存在密切的相互关系。
    \end{itemize}

    \textbf{D选项错误。}
    \begin{itemize}
        \item PAC层和支持层的提前还款风险是相互关联的，而不是独立的。换句话说，PAC层的表现会直接影响支持层的风险状况，二者之间的关系是动态的。
    \end{itemize}
\end{explanation}

\checklist{现代CMO结构中PAC层和支持层的关系。（记住PAC层和支持层的关系）}



\begin{question}
    A financial analyst at an investment firm is currently comparing the characteristics of Principal-Only (PO) securities and Interest-Only (IO) securities to provide valuable insights for investment decisions. Which of the following statements is correct?
\end{question}

\begin{choices}
    \item PO securities have larger cash flows at the beginning, which decrease over time.
    \item The total cash flow value received by IO securities throughout the life of the mortgage pool is always higher than initially expected.
    \item A higher prepayment rate will lead to a decrease in the yield of PO securities.
    \item Uncertainty in interest rates and prepayment rates typically results in a decrease in the value of IO securities and an increase in the value of PO securities.
\end{choices}

\begin{correctanswer}
    D
\end{correctanswer}

\begin{explanation}
    \textbf{A选项错误。}
    \begin{itemize}
        \item PO证券的现金流在初期较小，随着时间的推移而逐渐增加。这是因为在早期阶段，投资者回收本金的速度较慢，而随着时间的推移，回收速度会加快，从而导致现金流的增加。
    \end{itemize}

    \textbf{B选项错误。}
    \begin{itemize}
        \item IO证券在整个抵押贷款池的生命周期内收到的现金流价值可能会低于最初的预期，尤其是在提前还款率较高的情况下，投资者可能会面临现金流减少的风险。
    \end{itemize}

    \textbf{C选项错误。}
    \begin{itemize}
        \item 较高的提前还款率实际上会导致本金更快地被回收，从而产生更高的收益，而不是降低收益。对于PO证券来说，提前还款通常意味着投资者能够更快地收回本金，进而提高整体收益。
    \end{itemize}

    \textbf{D选项正确。}
    \begin{itemize}
        \item 利率和提前还款率的不确定性通常会导致IO证券的价值降低，因为投资者对未来现金流的预期变得不确定；而PO证券的价值则可能因本金回收的加速而增加，从而使其在不确定性中更具吸引力。
    \end{itemize}
\end{explanation}

\checklist{PO和IO证券的现金流特性。（记住现金流特性的差异）}


\begin{question}
    A financial analyst is evaluating the characteristics of Principal-Only (PO) securities and Interest-Only (IO) securities to determine their respective risks and returns. Which of the following statements is correct?
\end{question}

\begin{choices}
    \item Only IO securities carry the risk of loss.
    \item The cash flows of PO securities decrease over time.
    \item A higher prepayment rate will lead to higher yields for IO securities.
    \item Uncertainty in interest rates and prepayment rates leads to a decrease in the value of both PO and IO securities.
\end{choices}

\begin{correctanswer}
    C
\end{correctanswer}

\begin{explanation}
    \textbf{A选项错误。}
    \begin{itemize}
        \item IO证券在整个抵押贷款池的生命周期内收到的现金流价值可能低于最初预期，而PO证券的整个面值最终支付给投资者，因此PO证券也存在风险。
    \end{itemize}

    \textbf{B选项错误。}
    \begin{itemize}
        \item PO证券的现金流开始时较小，随着时间的推移而增加，因为随着抵押贷款支付中本金部分的增长而增加。
    \end{itemize}
    
    \textbf{C选项正确。}
    \begin{itemize}
        \item 较高的提前还款率导致本金更快回收，从而产生更高的收益，这适用于PO证券。
    \end{itemize}
    
    \textbf{D选项错误。}
    \begin{itemize}
        \item 利率和提前还款率的不确定性导致IO的价值降低，PO的价值增加，而不是两者都降低。
    \end{itemize}
\end{explanation}

\checklist{PO和IO证券的风险和回报特性。（记住现金流特性的差异）}


   


\begin{question}
    A financial analyst is examining the structure of the To Be Announced (TBA) market, which is a subset of the mortgage-backed securities (MBS) market. Which of the following statements most accurately characterizes the structure of the TBA market?
\end{question}

\begin{choices}
    \item TBAs are comprised of highly diverse securities thereby reducing the risk to investors.
    \item The composition of a TBA is announced at the time the trade bids on it.
    \item The TBA market is essentially a forward market in which buyer and seller agree to exchange a pool of securities for cash in the future.
    \item There are actually two separate TBA markets, one for fixed-rate agency pools and the other for conventional adjustable-rate mortgage pools.
\end{choices}

\begin{correctanswer}
    C
\end{correctanswer}

\begin{explanation}
    \textbf{A选项错误。}
    \begin{itemize}
        \item TBAs并不是由高度多样化的证券组成。相反，TBA市场主要涉及特定类型的抵押贷款支持证券，因此不能有效降低投资者的风险。投资者在TBA市场中面临的风险主要与特定证券的信用质量和市场条件相关。
    \end{itemize}

    \textbf{B选项错误。}
    \begin{itemize}
        \item TBA的组成在交易时并不会被宣布。实际上，交易双方在达成协议时会确定具体的证券池，而不是在交易时公开。这种机制使得交易双方在交易前对证券的具体组成有一定的灵活性。
    \end{itemize}

    \textbf{C选项正确。}
    \begin{itemize}
        \item TBA市场本质上是一个远期市场，买卖双方同意在未来交换一组证券以换取现金。这种结构允许投资者在未来的某个时间点以约定的价格进行交易，从而提供了流动性和价格确定性。
    \end{itemize}

    \textbf{D选项错误。}
    \begin{itemize}
        \item TBA市场实际上是一个统一的市场，并没有两个独立的市场来处理固定利率和可调利率的抵押贷款池。所有的TBA交易都在同一个市场中进行，尽管可能会涉及不同类型的抵押贷款，但它们并不被视为两个独立的市场。
    \end{itemize}
\end{explanation}

\checklist{TBA市场的结构。（记住TBA市场的结构）}



\begin{question}
    An investor in MBS is evaluating two methods of financing a recent purchase of an MBS pool. One method involves entering into a repurchase agreement (repo) and selling the pool now. The other method involves purchasing a dollar roll. One advantage of purchasing a roll instead of entering into a repo agreement is that roll buyers:
\end{question}

\begin{choices}
    \item Receive the same securities back when the transaction is concluded.
    \item Receive the interest and principal repayments on the pool during the life of the roll.
    \item Are more similar to a secured borrower than a repo participant.
    \item Are subject to less risk of prepayment than a repo buyer.
\end{choices}

\begin{correctanswer}
    D
\end{correctanswer}

\begin{explanation}
    \textbf{A选项错误。}
    \begin{itemize}
        \item 在美元滚动交易中，交易双方并不要求返还相同的证券，而是可以交付具有相似特征的MBS证券。这种灵活性实际上是美元滚动交易的一个特点，而不是优势。
    \end{itemize}

    \textbf{B选项错误。}
    \begin{itemize}
        \item 在美元滚动交易期间，投资者放弃了本金和利息支付权。虽然这看似是一个劣势，但这种结构实际上帮助投资者降低了提前还款风险，并且通常能够通过较低的回购价格获得补偿。
    \end{itemize}

    \textbf{C选项错误。}
    \begin{itemize}
        \item 美元滚动交易在本质上是一种特殊的回购协议，它提供了一种以MBS为抵押品的短期融资方式。这种结构与担保借款有所不同，因为它涉及证券的实际买卖。
    \end{itemize}

    \textbf{D选项正确。}
    \begin{itemize}
        \item 美元滚动交易中，买家确实面临较低的提前还款风险，这是因为：
        \begin{itemize}
            \item 在滚动期间放弃本金和利息支付，避免了这期间的提前还款风险
            \item 能够以低于销售价格的价格回购抵押贷款，提供了额外的价格保护
            \item 交易结构允许在回购时选择不同的证券，增加了风险管理的灵活性
        \end{itemize}
    \end{itemize}
\end{explanation}

\checklist{美元滚动交易的优势。（美元滚动交易中，买家确实面临较低的提前还款风险）}











\fi




\iftrue
\section{MBS考点3:Valuation of MBSs pool}
\setcounter{question}{0}



\begin{question}
    Emma, a structured products analyst at a global investment bank, is evaluating a new residential mortgage-backed security (RMBS) with a pool size of 800 million dollars. Her team needs to assess the prepayment risk characteristics of the pool. Which of the following statements about MBS analysis is correct?
    \end{question}
    
    \begin{choices}
        \item For Ginnie Mae MBS valuation, both prepayment risk and credit risk modeling should be prioritized as they have equal importance in price determination.
        
        \item The weighted-average FICO score of the mortgage pool primarily reflects the credit quality of the lending institutions in the securitization.
        
        \item When Single Monthly Mortality (SMM) rate increases, IO strip values rise due to increased interest payments, while PO strip values fall due to faster principal repayment.
        
        \item In prepayment modeling, FICO scores and LTV ratios are key inputs for default probability estimation, while loan seasoning helps predict voluntary prepayment patterns.
    \end{choices}
    
    \begin{correctanswer}
    D
    \end{correctanswer}
    
    \begin{explanation}
        \textbf{A选项错误。}
        \begin{itemize}
            \item Agency MBS (如Ginnie Mae发行的) 具有政府担保，信用风险几乎为零。估值主要考虑提前还款风险：
            \[Price_{\text{MBS}} = E\left[\sum_{t=1}^T \frac{CF_t}{(1+r)^t}\right]\]
            其中$CF_t$受提前还款行为影响，而不考虑违约风险。
        \end{itemize}
        
        \textbf{B选项错误。}
        \begin{itemize}
            \item FICO分数反映借款人信用质量，计算公式为：
            \[\text{FICO}_{\text{pool}} = \sum_{i=1}^n w_i \times \text{FICO}_i\]
            其中$w_i$为第$i$笔贷款的权重，$\text{FICO}_i$为对应的借款人信用分数。
        \end{itemize}
        
        \textbf{C选项错误。}
        \begin{itemize}
            \item SMM与证券价值的关系：
            \[\text{IO}_{\text{value}} \propto \sum_{t=1}^T r \times (1-\text{CPR})^t \times \text{Balance}_t\]
            \[\text{PO}_{\text{value}} \propto \sum_{t=1}^T \text{CPR} \times \text{Balance}_t\]
            其中$\text{CPR} = 1-(1-\text{SMM})^{12}$，$\propto$表示“正比于”。SMM上升导致IO价值下降，PO价值上升。
        \end{itemize}
        
        \textbf{D选项正确。}
        \begin{itemize}
            \item 违约概率(PD)和提前还款率(CPR)的建模关系：
            \[\text{PD} = f(\text{FICO}, \text{LTV}, \text{其他因素})\]
            \[\text{CPR} = g(\text{贷款年限}, \text{市场利率}, \text{季节性因素})\]
            这些变量共同决定了MBS的现金流特征。
        \end{itemize}
        
        \end{explanation}
    
    \checklist{MBS的提前还款风险建模。（了解FICO分数、LTV比率、贷款期限，对违约风险和提前还款行为的影响机制。）}



    \begin{question}
        A portfolio manager at a quantitative hedge fund is evaluating an agency MBS with a pool size of \$500 million. She needs to identify factors that could increase prepayment risk and affect the portfolio's cash flows. Which of the following scenarios would most likely increase prepayment rates?
    \end{question}
        
        \begin{choices}
            \item The weighted-average FICO score of borrowers in the mortgage pool declines from 760 to 690 due to increased credit card utilization.
            
            \item The mortgage pool consists mainly of 2-year seasoned loans with an average loan-to-value ratio of 85\%.
            
            \item The 90-day delinquency rate in comparable mortgage pools rises from 0.4\% to 1.1\% amid economic uncertainty.
            
            \item The Federal Reserve implements a dovish policy shift, causing the 30-year mortgage rate to decrease by 125 basis points.
        \end{choices}
        
        \begin{correctanswer}
        D
        \end{correctanswer}
        
        \begin{explanation}
       
        
        \textbf{A选项错误。}
        \begin{itemize}
            \item FICO分数从760降至690意味着借款人信用质量显著下降，这会导致再融资难度增加，新贷款利率可能上浮75-100个基点，实际上会降低而不是增加提前还款的可能性。
        \end{itemize}
        
        \textbf{B选项错误。}
        \begin{itemize}
            \item 2年期的贷款和85\%的高贷款价值比表明借款人积累的房屋净值有限，大部分早期还款用于支付利息，这种情况下借款人缺乏再融资的条件和动力，不会增加提前还款。
        \end{itemize}
        
        \textbf{C选项错误。}
        \begin{itemize}
            \item 违约率从0.4\%上升到1.1\%虽然显著，但对机构MBS影响有限，因为政府担保机构会回购违约贷款，且违约赎回与自愿提前还款是不同的现象，实际上可能降低正常的提前还款活动。
        \end{itemize}
        
        \textbf{D选项正确。}
        \begin{itemize}
            \item 125个基点的利率下降将显著降低借款人的月供支出（约12-15\%），即使考虑交易成本，也为借款人提供了明显的再融资收益，这是刺激提前还款的最主要因素。
        \end{itemize}
        
        \end{explanation}
        
        \checklist{分析MBS提前还款风险的关键因素。（弄明白再融资和提前还款的逻辑关系）}




        \begin{question}
            In the context of mortgage-backed securities (MBS), modeling prepayment behavior involves complex mathematical relationships that depend not only on interest rates but also on various factors. A commonly used prepayment rate model is:
            \[
            \text{Prepayment Rate} = \frac{1}{a + b \cdot e^{-c \cdot T}}
            \]
            where \( a, b, c \) are parameters determined from empirical data, and \( T \) is time. The prepayment rate increases as interest rates decline but tends to stabilize after reaching a certain level. Please explain why the prepayment rate increases initially with declining interest rates but ultimately approaches a stable value.
        \end{question}
        
        \begin{choices}
            \item In the early stages of declining interest rates, borrowers are inclined to refinance at lower rates, thus increasing the prepayment rate. However, over time, most eligible borrowers have already refinanced, leading to a stabilization of the prepayment rate.
            \item In the early stages of declining interest rates, borrowers have no incentive to prepay because the new loan rates may not be significantly lower than their current loan rates, so the prepayment rate does not increase.
            \item The increase in the prepayment rate is entirely driven by market sentiment and is unrelated to changes in interest rates, meaning that even if rates decline, borrowers may choose to prepay for other reasons.
            \item The prepayment rate will never stabilize; instead, it will continue to increase as interest rates decline, as borrowers always tend to refinance at lower rates.
        \end{choices}
        
        \begin{correctanswer}
            A
        \end{correctanswer}
        
        \begin{explanation}
            \textbf{A选项正确。}
            \begin{itemize}
                \item 在利率下降初期，借款人倾向于利用更低的利率重新融资，从而增加提前还款率，但随着时间的推移，大多数符合条件的借款人已经重新融资，导致提前还款率趋于稳定。
            \end{itemize}
        
            \textbf{B选项错误。}
            \begin{itemize}
                \item 实际上，利率下降初期，借款人有动力提前还款以利用更低的利率重新融资，因此提前还款率会增加。
            \end{itemize}
        
            \textbf{C选项错误。}
            \begin{itemize}
                \item 虽然市场情绪可能影响提前还款率，但利率变化是提前还款率增加的一个重要因素，特别是在利率下降时。
            \end{itemize}
        
            \textbf{D选项错误。}
            \begin{itemize}
                \item 根据模型，提前还款率在达到一定水平后会趋于平稳，不会持续无限增加，因为大多数符合条件的借款人已经重新融资。
            \end{itemize}
        \end{explanation}
        
        \checklist{提前还款率模型。（利率下降初期增加提前还款率。随着时间的推移提前还款率趋于稳定。）}


        \begin{question}
            A risk analyst at a mortgage investment firm is studying various methods to address prepayment path dependency in mortgage-backed securities. Which of the following methods is often used to avoid the problems associated with prepayment path dependency?
        \end{question}
        
        \begin{choices}
            \item Monte Carlo simulation.
            \item Error-correction model tree.
            \item Cox-Ingersoll-Ross tree design.
            \item Bernard and Schwartz simulation.
        \end{choices}
        
        \begin{correctanswer}
            A
        \end{correctanswer}
        
        \begin{explanation}
            \textbf{A选项正确。}
            \begin{itemize}
                \item 蒙特卡洛模拟是处理提前还款路径依赖性问题的有效工具，因为它可以模拟大量可能的利率路径和提前还款情景，并能够捕捉复杂的路径依赖特征，同时考虑多个影响因素的相互作用。
            \end{itemize}
        
            \textbf{B选项错误。}
            \begin{itemize}
                \item 误差修正模型树虽然在某些金融建模中有应用，但并不特别适合处理提前还款路径依赖性问题。
            \end{itemize}
        
            \textbf{C选项错误。}
            \begin{itemize}
                \item Cox-Ingersoll-Ross树状模型主要用于利率建模，虽然与抵押贷款估值相关，但不是专门用来解决提前还款路径依赖性问题的工具。
            \end{itemize}
        
            \textbf{D选项错误。}
            \begin{itemize}
                \item Bernard和Schwartz模拟虽然在金融领域有其应用，但在处理提前还款路径依赖性问题上不如蒙特卡洛模拟普遍和有效。
            \end{itemize}
        \end{explanation}

        \checklist{处理MBS提前还款路径依赖性问题的方法。（蒙特卡洛模拟是处理提前还款路径依赖性问题的有效工具。）}


        \begin{question}
            A mortgage analyst at a financial institution is studying various factors affecting prepayment behavior in mortgage-backed securities (MBS). Which of the following statements correctly describes a characteristic of prepayment behavior?
        \end{question}
        
        \begin{choices}
            \item The likelihood of refinancing is positively correlated with market interest rate increases and negatively correlated with improvements in borrower credit quality.
            \item The turnover rate is not affected by seasonality and remains stable throughout the year.
            \item Default is considered a form of prepayment and affects MBS cash flows.
            \item Curtailment typically occurs at the beginning of the loan term to reduce the loan balance.
        \end{choices}
        
        \begin{correctanswer}
            C
        \end{correctanswer}
        
        \begin{explanation}
            \textbf{A选项错误。}
            \begin{itemize}
                \item 再融资的可能性与市场利率的下降（而不是上升）成正比，与借款人信用质量的提高成正比。这是因为较低的市场利率和较好的信用质量都会增加借款人获得更优惠贷款条件的机会。
            \end{itemize}
        
            \textbf{B选项错误。}
            \begin{itemize}
                \item 周转率实际上具有明显的季节性特征，通常在夏季高于冬季。这主要是因为房屋买卖活动在夏季更为活跃。
            \end{itemize}
        
            \textbf{C选项正确。}
            \begin{itemize}
                \item 违约确实被视为提前还款的一种形式，因为它会导致贷款提前终止，从而影响MBS的现金流。这是MBS风险评估中的一个重要考虑因素。
            \end{itemize}
        
            \textbf{D选项错误。}
            \begin{itemize}
                \item 削减（部分提前还款）通常发生在贷款期限的最后几年，而不是开始阶段。这种行为可能导致高达35\%的提前还款率。
            \end{itemize}
        \end{explanation}
\checklist{MBS提前还款的因素。（违约确实被视为提前还款的一种形式，因为它会导致贷款提前终止，从而影响MBS的现金流。）}


\begin{question}
    A fixed-income analyst at an investment bank is studying the Option-Adjusted Spread (OAS) methodology for MBS valuation. Which of the following statements MOST LIKELY misrepresents the characteristics of  the OAS valuation procedure?
\end{question}

\begin{choices}
    \item The OAS valuation procedure involves making an initial estimate of the OAS, followed by conducting a Monte Carlo simulation using the yield on US Treasury bills as the discount rate.
    \item OAS represents the expected yield spread of fixed-income securities over the risk-free rate and is used to measure the relative valuation of MBS.
    \item The mean reversion characteristic of OAS implies that OAS estimates are repeatedly adjusted until the simulated price equals the market price.
    \item A limitation of the OAS valuation method is that OAS itself implies an assumption that the spread remains constant over time, although OAS actually changes over time.
\end{choices}

\begin{correctanswer}
    A
\end{correctanswer}

\begin{explanation}
    \textbf{A选项错误。}
    \begin{itemize}
        \item OAS估值程序中，在进行蒙特卡洛模拟时，正确的贴现率应该是国债收益率加上初始OAS估计值，而不是仅使用国债收益率。这个差异很重要，因为OAS反映了证券相对于无风险利率的额外收益，忽略这部分将导致估值偏差。
    \end{itemize}

    \textbf{B选项正确，不选。}
    \begin{itemize}
        \item 这个陈述准确定义了OAS的本质：它是固定收益证券预期收益率超过无风险收益率的差额。这个指标确实被广泛用于衡量MBS的相对估值，因为它考虑了期权特征对价值的影响。
    \end{itemize}

    \textbf{C选项正确，不选。}
    \begin{itemize}
        \item OAS估值过程确实具有均值复归特性，需要通过反复迭代调整OAS估计值。这个过程会持续到模拟得到的价格与市场观察到的价格相等为止，这是OAS估值的一个重要特征。
    \end{itemize}

    \textbf{D选项正确，不选。}
    \begin{itemize}
        \item 这个选项准确指出了OAS方法的一个重要局限性：它假设利差在一段时间内保持不变。实际上，OAS会随时间变化，这种假设可能导致估值偏差，这是使用OAS方法时需要注意的重要问题。
    \end{itemize}
\end{explanation}

\checklist{OAS估值程序。（OAS估值程序中，在进行蒙特卡洛模拟时，正确的贴现率应该是国债收益率加上初始OAS估计值，而不是仅使用国债收益率。）}





\begin{question}
    Emily, a fixed-income analyst, recently completed a Monte Carlo simulation analysis of a CMO tranche. Her analysis includes three equally weighted paths, with the present value of each calculated using different spreads. If the actual market price of the CMO tranche is 70.00, what is the option-adjusted spread (OAS)?
    \begin{center}
        \begin{tabular}{|c|c|c|c|}
        \hline
        \textbf{Path} & \textbf{PV at 55 bps} & \textbf{PV at 65 bps} & \textbf{PV at 75 bps} \\
        \hline
        1 & 72 & 70 & 68 \\
        2 & 74 & 72 & 70 \\
        3 & 70 & 68 & 66 \\
        \hline
        \end{tabular}
        \end{center}
    \end{question}
    
    \begin{choices}
        \item 55 basis points
        \item 65 basis points
        \item 75 basis points
        \item 80 basis points
    \end{choices}
    
    
    
    \begin{correctanswer}
    B
    \end{correctanswer}
    
 \begin{explanation}
    \textbf{步骤1：理解OAS的概念。}
    \begin{itemize}
        \item OAS是使模拟价格等于市场价格的利差
        \item 在本题中，需要找到使蒙特卡洛模拟的平均现值等于市场价格70.00的利差
    \end{itemize}

    \textbf{步骤2：分析三个利差水平下的平均现值。}
    \begin{itemize}
        \item 在55 bps时：\[\frac{72 + 74 + 70}{3} = 72.00 > 70.00\]
        \item 在65 bps时：\[\frac{70 + 72 + 68}{3} = 70.00 = 70.00\]
        \item 在75 bps时：\[\frac{68 + 70 + 66}{3} = 68.00 < 70.00\]
    \end{itemize}

        \textbf{步骤3：确定OAS。}
    \begin{itemize}
        \item 观察到65 bps时的平均现值恰好等于市场价格70.00
        \item 这说明65 bps就是我们要找的OAS
               \item 验证价格回归过程：
            \begin{itemize}
                \item 利差与价格呈反向关系：利差↑ → 价格↓
                \item 通过三个利差水平的测试：
                    \begin{itemize}
                        \item 55 bps：价格72.00 > 目标70.00，利差需要提高
                        \item 65 bps：价格70.00 = 目标70.00，找到正确利差
                        \item 75 bps：价格68.00 < 目标70.00，利差过高
                    \end{itemize}
                \item 结论：65 bps是使模拟价格等于市场价格的OAS
            \end{itemize}
    \end{itemize}
\end{explanation}
    
    \checklist{计算OAS的步骤。（理解OAS的原理）}



    \begin{question}
        A mortgage analyst is evaluating factors that could influence prepayment rates in a portfolio of Mortgage-Backed Securities (MBS). Which of the following factors is most likely to lead to increased prepayments in Mortgage-Backed Securities (MBS)?
    \end{question}
        
        \begin{choices}
            \item A stable housing market encourages borrowers to maintain their current mortgages.
            \item Volatile housing prices prompt borrowers to refinance or sell their properties.
            \item Rising interest rates discourage refinancing activities.
            \item Approaching maturity of the MBS has minimal impact on prepayment behavior.
        \end{choices}
        
        \begin{correctanswer}
        B
        \end{correctanswer}
        
        \begin{explanation}
        
        \textbf{A选项错误。}
        \begin{itemize}
            \item 稳定的住房市场通常不会激励借款人进行再融资或提前还款，因为缺乏价格波动带来的机会。
        \end{itemize}
        
        \textbf{B选项正确。}
        \begin{itemize}
            \item 住房价格的剧烈波动可能促使借款人再融资以利用更高的房屋价值，或出售房产以获利，从而增加MBS的提前还款。
        \end{itemize}
        
        \textbf{C选项错误。}
        \begin{itemize}
            \item 利率上升通常会减少借款人再融资的动机，因为他们更倾向于保持现有的较低利率贷款。
        \end{itemize}
        
        \textbf{D选项错误。}
        \begin{itemize}
            \item 接近MBS的到期日可能会影响市场情绪，但通常不会直接导致提前还款增加。
        \end{itemize}
        
        \end{explanation}
        
        \checklist{分析影响MBS提前还款的因素（包括市场条件和利率变化）}
            



\begin{question}
    A quantitative analyst at a mortgage investment firm is tasked with improving their MBS valuation models. During their research, they discover that traditional valuation methods may not fully capture the complexity of these securities. In their analysis of different methodologies, they focus particularly on path dependency characteristics. Which of their findings about path dependency and valuation methods is correct?
\end{question}

\begin{choices}
    \item Monte Carlo simulation provides limited value in path dependency analysis due to its focus on single-period scenarios.
    \item The burnout effect indicates that prepayment rates tend to decrease when refinancing rates drop substantially, as borrowers who are likely to refinance have already done so in previous rate cycles.
    \item Path dependency primarily affects short-term pricing adjustments in MBS valuation, with minimal impact on long-term value.
    \item Monte Carlo simulation's primary advantage lies in its computational speed rather than its ability to capture path dependency characteristics.
\end{choices}

\begin{correctanswer}
    B
\end{correctanswer}

\begin{explanation}
    \textbf{A选项错误。}
    \begin{itemize}
        \item 蒙特卡洛模拟实际上非常适合处理路径依赖问题，因为它可以模拟多期情景，而不是仅限于单期情景。这种方法能够有效捕捉MBS估值中的路径依赖特征。
    \end{itemize}

    \textbf{B选项正确。}
    \begin{itemize}
        \item 燃尽效应（burnout effect）确实表明，当再融资利率显著下降时，提前还款率往往不会大幅上升。这是因为在之前的低利率周期中，倾向于再融资的借款人已经完成了再融资。剩余的借款人对利率变化的敏感度较低，即使利率进一步下降，也不太可能进行再融资。
    \end{itemize}

    \textbf{C选项错误。}
    \begin{itemize}
        \item 路径依赖对MBS估值的影响是长期的和根本的，而不仅仅影响短期定价调整。利率变化的路径会显著影响MBS的整体价值和现金流特征。
    \end{itemize}

    \textbf{D选项错误。}
    \begin{itemize}
        \item 蒙特卡洛模拟的主要优势在于其能够有效捕捉路径依赖特征，而不是计算速度。实际上，为了准确模拟路径依赖性，蒙特卡洛模拟往往需要大量的计算时间。
    \end{itemize}
\end{explanation}

\checklist{MBS估值中的路径依赖问题。（燃尽效应的特征）}


\begin{question}
    A senior analyst at a fixed-income investment firm has been tasked with training new team members about MBS valuation methods. During a presentation on Option-Adjusted Spread (OAS) analysis, several team members raise questions about the methodology's limitations. The analyst reviews their notes and needs to identify which statement about OAS characteristics requires correction. Which of the following statements MOST LIKELY misrepresents the characteristics of OAS?
\end{question}

\begin{choices}
    \item OAS is based on Monte Carlo simulation results, which may inherently contain model risk.
    \item The adjustment of interest rate paths may be subjective and affected by model risk.
    \item OAS implies an assumption that spreads remain constant over a period of time, although OAS actually changes over time.
    \item Prepayment rates are closely related to OAS modeling, and MBS prepayment rates can be accurately estimated.
\end{choices}

\begin{correctanswer}
    D
\end{correctanswer}

\begin{explanation}
    \textbf{A选项正确，不选。}
    \begin{itemize}
        \item OAS确实是基于蒙特卡洛模拟计算得出的，这种方法本身就可能存在模型风险。模型风险主要来源于模拟过程中使用的各种假设可能与实际情况不符。
    \end{itemize}

    \textbf{B选项正确，不选。}
    \begin{itemize}
        \item 在OAS分析中，利率路径的调整确实包含主观判断的成分。这种主观性可能导致不同分析师得出不同的结果，增加了模型风险。
    \end{itemize}

    \textbf{C选项正确，不选。}
    \begin{itemize}
        \item OAS分析确实假设在某一时期内利差保持不变。但实际上，市场环境的变化会导致OAS随时间发生变化，这是OAS分析的一个重要局限。
    \end{itemize}

    \textbf{D选项错误。}
    \begin{itemize}
        \item 提前还款率确实与OAS建模密切相关，但说它可以被准确估计是错误的。提前还款行为受多种因素影响，如借款人行为、经济环境、市场条件等，这些因素的复杂性使得准确预测变得极其困难。提前还款率的不确定性是MBS估值中的一个主要风险因素。
    \end{itemize}
\end{explanation}

\checklist{OAS的不足。（OAS方法有模型风险）}











\end{document}